%!TeX root = ../principal.tex
%!TeX encoding = utf-8
\chapter{Modelos}
\label{cap3}
Comparamos três famílias de modelos representativas, que incorporam diferentes vieses indutivos:

\begin{itemize}

    \item \textbf{Redes de Cápsulas (Efficient-CapsNet)} \citar{mazzia2021efficient}, projetadas com fortes vieses de equivariância e de hierarquia parte-todo \citar{sabour2017dynamic}, e historicamente apontadas como vantajosas em cenários com poucos dados e na robustez a transformações.

    \item \textbf{CNNs Residuais (ResNet-18)} \citar{he2016resnet}, uma das arquiteturas convolucionais modernas mais amplamente utilizadas, representando o viés indutivo dominante de equivariância à translação e aprendizado hierárquico de características por meio de conexões residuais.

    \item \textbf{Transformadores de Visão (DeiT-Tiny)} \citar{dosovitskiy2021vit,touvron2021deit}, arquiteturas com vieses indutivos incorporados mínimos, que dependem principalmente dos dados e de mecanismos de atenção.
\end{itemize}

As três arquiteturas cabem em uma única GPU de 6\,GB e são usadas na
configuração padrão de cada família, sem pareamento de capacidade. O desnível
que resulta é grande: a ResNet-18 tem cerca de $2{,}8$ vezes os parâmetros do
Efficient-CapsNet, e o DeiT-Tiny cerca de $1{,}4$ vez
(Tabela~\ref{tab:model-complexity}). Considerando apenas o extrator e as
cápsulas, sem o decodificador de reconstrução, o modelo capsular tem cerca de
$0{,}16$~M de parâmetros, dezenas de vezes menos que as linhas de base. As
diferenças de desempenho, portanto, refletem o viés indutivo e a capacidade de
cada arquitetura, sem que seja possível separar os dois.

% ---------------------------------------------------------------------
\section{Arquitetura e Modificações}
\label{sec:protocol:models}

As três arquiteturas são disponibilizadas por meio de uma única interface
\texttt{forward(x,\allowbreak{} y\_true,\allowbreak{} mode)}, de modo que o treinador, a infraestrutura
de registro e o sistema de avaliação sejam compartilhados integralmente entre
os modelos; apenas a arquitetura e a função de perda diferem.

\subsection{Efficient-CapsNet}
\label{sec:protocol:models:ecaps}

A rede capsular adotada segue a arquitetura Efficient-CapsNet
\citar{mazzia2021efficient}. O extrator convolucional é composto por
quatro camadas, todas seguidas de normalização em lote e ativação
ReLU: uma convolução $5\times5$ com 32 canais, duas convoluções
$3\times3$ com 64 canais e uma convolução $3\times3$ com 128 canais e
passo 2. Nenhuma delas utiliza preenchimento, de modo que o mapa de
características é progressivamente reduzido.

A camada \textit{PrimaryCaps} converte esse mapa em cápsulas por meio
de uma convolução \textit{depthwise} cujo núcleo tem exatamente as
dimensões espaciais remanescentes do extrator, colapsando o mapa para
$1\times1$; a saída é reorganizada em 16 cápsulas de dimensão 8 e
normalizada pela função \textit{squash}. Em seguida, a camada
\textit{RoutingCaps} substitui o roteamento iterativo por acordo por um
mecanismo de autoatenção não iterativo: os vetores de predição são
obtidos por um produto tensorial entre as cápsulas de entrada e a
matriz de pesos $\mathbf{W} \in \mathbb{R}^{N_1 \times N_0 \times D_0
\times D_1}$, e os coeficientes de acoplamento resultam da
autossimilaridade entre esses vetores, escalonada por $\sqrt{D_1}$,
seguida de \textit{softmax} e de um termo de viés aprendido. A saída
consiste em uma cápsula de dimensão 16 por classe. Na configuração do
CIFAR-10 há ainda uma cápsula de fundo, que nunca é alvo do treino e é
descartada na classificação.
% VERIFICAR: use_background_class=True nos YAML do CIFAR-10 e False (ou
% ausente) nos do D-Fire. A contagem de parâmetros da
% Tabela tab:model-complexity (3,94 M e 3,85 M) só fecha com essa combinação.

O comprimento euclidiano de cada cápsula de saída é utilizado como
escore da classe correspondente. Adicionalmente, um decodificador
totalmente conectado (camadas de 512 e 1024 unidades, seguidas de uma
projeção de $3\times32\times32$ saídas com ativação sigmoide) reconstrói a
imagem a partir das cápsulas mascaradas pelo rótulo verdadeiro durante o
treinamento, ou pela cápsula de maior comprimento durante a inferência. A
perda de reconstrução atua como regularizador, conforme descrito na
Seção~\ref{sec:protocol:training}. O decodificador concentra cerca de
$3{,}7$~M dos parâmetros do modelo, mais de $95\%$ do total.

Para entradas maiores que $32\times32$ foram necessárias duas adaptações.
A primeira está no extrator. Sem ela, a única redução espacial seria o passo 2
da última convolução, o mapa final teria $107\times107$ posições e o custo
cresceria com o quadrado da resolução. Por isso inserimos um
\textit{max-pooling} $2\times2$ após cada uma das três primeiras convoluções
quando a entrada é de $224\times224$. O mapa final volta a ter $12\times12$
posições, próximo das $11\times11$ da configuração original em $32\times32$,
que permanece inalterada. O núcleo da convolução \textit{depthwise} das
cápsulas primárias acompanha esse tamanho: $11\times11$ no CIFAR-10 e
$12\times12$ no D-Fire.

A segunda adaptação está no decodificador. Reconstruir diretamente em
$3\times224\times224$ exigiria uma última camada com cerca de 154~M de
parâmetros, dezenas de vezes o resto do modelo. No D-Fire, o decodificador
reconstrói em $32\times32$, e a saída é ampliada por interpolação bilinear
para $224\times224$ antes do cálculo da perda de reconstrução. O decodificador
tem, assim, praticamente o mesmo tamanho nos dois conjuntos.
% VERIFICAR: recon_size: 32 nos YAML do D-Fire (é o que a contagem de
% 3,85 M parâmetros indica).

Mesmo com essas adaptações, o custo capsular não escala como o das linhas de
base: o número de operações do Efficient-CapsNet cresce cerca de nove vezes
entre os dois conjuntos, contra três a quatro vezes para a ResNet-18 e o
DeiT-Tiny (Tabela~\ref{tab:model-complexity}). A diferença é considerada na
Seção~\ref{sec:protocol:threats}.

\subsection{ResNet-18}
\label{sec:protocol:models:resnet}

A segunda arquitetura de referência é a ResNet-18 \citar{he2016resnet},
composta por quatro estágios de blocos residuais básicos com dois
caminhos convolucionais $3\times3$ cada, conexões de atalho por
identidade e normalização em lote.

Para as entradas de baixa resolução do CIFAR-10, adota-se a adaptação
usual do extrator inicial: a convolução $7\times7$ com passo 2 é
substituída por uma convolução $3\times3$ com passo 1 e preenchimento
1, e a camada de \textit{max-pooling} subsequente é removida. Sem essa
alteração, a imagem $32\times32$ seria reduzida a $8\times8$ antes do
primeiro bloco residual, descartando informação espacial essencial. Os
pesos da nova convolução são inicializados segundo a distribuição
normal de He. No D-Fire, com entradas de $224\times224$, a resolução para
a qual a arquitetura foi projetada, o extrator original é mantido.

\subsection{DeiT-Tiny}
\label{sec:protocol:models:deit}

A terceira arquitetura é a DeiT-Tiny \citar{touvron2021deit}, um
\textit{Vision Transformer} com 12 blocos, dimensão de imersão 192 e 3
cabeças de atenção. A destilação por token proposta no trabalho
original não é empregada aqui, uma vez que exigiria um modelo professor
e comprometeria a comparação entre arquiteturas treinadas do zero.

O tamanho do \textit{patch} é ajustado à resolução de cada conjunto de
dados: \textit{patches} de $4\times4$ para as entradas $32\times32$ do
CIFAR-10, resultando em 64 tokens, e \textit{patches} de $16\times16$
para as entradas $224\times224$ do D-Fire, resultando em 196 tokens.
Como o número de tokens do CIFAR-10 difere da configuração original da
arquitetura, as imersões posicionais são reinstanciadas com as
dimensões correspondentes e aprendidas do zero. Em ambos os conjuntos o
modelo parte de inicialização aleatória, coerente com o protocolo de
treinamento adotado para todos os modelos.